\chapter{2011年大纲卷（理）}

\section{单选题}

\begin{problem}
复数 \(z = 1 + \i\)，\(\overline{z}\) 为 \(z\) 的共轭复数. 则 \(z\overline{z} - z - 1 =\)（\quad）

\choices
    {\(-2\i\)}
    {\(-\i\)}
    {\(\i\)}
    {\(2\i\)}
\end{problem}

\begin{answer}
B
\end{answer}


\begin{problem}
函数 \( y = 2\sqrt{x} \) （\(x \geqslant 0\)）的反函数为（\quad）

\choices
    {\( y = \frac{x^2}{4} \)（\(x \in \R\)）}
    {\( y = \frac{x^2}{4} \)（\(x \geqslant 0\)）}
    {\( y = 4x^{2} \)（\(x \in \R\)）}
    {\( y = 4x^{2} \)（\(x \geqslant 0\)）}
\end{problem}

\begin{answer}
B
\end{answer}


\begin{problem}
下面四个条件中，使 \(a > b\) 成立的充分而不必要的条件是（\quad）

\choices
    {\(a > b + 1\)}
    {\(a > b - 1\)}
    {\(a^2 > b^2\)}
    {\(a^3 > b^3\)}
\end{problem}

\begin{answer}
A
\end{answer}


\begin{problem}
设 \(S_{n}\) 为等差数列 \(\{a_{n}\}\) 的前 \(n\) 项和，若 \(a_{1} = 1\)，公差 \(d = 2\)，\(S_{k+2} - S_{k} = 24\)，则 \(k =\)（\quad）

\choices
    {\(8\)}
    {\(7\)}
    {\(6\)}
    {\(5\)}
\end{problem}

\begin{answer}
D
\end{answer}


\begin{problem}
设函数 \( f(x) = \cos \omega x  \) （\(\omega > 0\)），将 \( y = f(x) \) 的图象向右平移 \( \frac{\pi}{3} \) 个单位长度后，所得的图象与原图象重合，则 \( \omega \) 的最小值等于（\quad）

\choices
    {\( \frac{1}{3} \)}
    {\(3\)}
    {\(6\)}
    {\(9\)}
\end{problem}

\begin{answer}
C
\end{answer}


\begin{problem}
已知直二面角 \(\alpha - l - \beta\)，点 \(A \in \alpha\)，\(AC \perp l\)，\(C\) 为垂足，\(B \in \beta\)，\(BD \perp l\)，\(D\) 为垂足. 若 \(AB = 2\)，\(AC = BD = 1\)，则 \(D\) 到平面 \(ABC\) 的距离等于（\quad）

\choices
    {\(\frac{\sqrt{2}}{3}\)}
    {\(\frac{\sqrt{3}}{3}\)}
    {\(\frac{\sqrt{6}}{3}\)}
    {\(1\)}
\end{problem}

\begin{answer}
C
\end{answer}


\begin{problem}
某同学有同样的画册 \(2\) 本，同样的集邮册 \(3\) 本，从中取出 \(4\) 本赠送给 \(4\) 位朋友，每位朋友 \(1\) 本，则不同的赠送方法共有（\quad）

\choices
    {\(4\) 种}
    {\(10\) 种}
    {\(18\) 种}
    {\(20\) 种}
\end{problem}

\begin{answer}
B
\end{answer}


\begin{problem}
曲线 \( y = \e^{-2x} + 1 \) 在点 \((0,2)\) 处的切线与直线 \( y = 0 \) 和 \( y = x \) 围成的三角形的面积为（\quad）

\choices
    {\(\frac{1}{3}\)}
    {\(\frac{1}{2}\)}
    {\(\frac{2}{3}\)}
    {\(1\)}
\end{problem}

\begin{answer}
A
\end{answer}


\begin{problem}
设 \( f(x) \) 是周期为 \(2\) 的奇函数，当 \( 0 \leqslant x \leqslant 1 \) 时，\( f(x) = 2x(1 - x) \)，则 \( f\left(-\frac{5}{2}\right) = \)（\quad）

\choices
    {\(-\frac{1}{2}\)}
    {\(-\frac{1}{4}\)}
    {\(\frac{1}{4}\)}
    {\(\frac{1}{2}\)}
\end{problem}

\begin{answer}
A
\end{answer}


\begin{problem}
已知抛物线 \(C: y^{2} = 4x\) 的焦点为 \(F\)，直线 \(y = 2x - 4\) 与 \(C\) 交于 \(A\)，\(B\) 两点，则 \(\cos \angle AFB =\)（\quad）

\choices
    {\(\frac{1}{5}\)}
    {\(\frac{3}{5}\)}
    {\(-\frac{3}{5}\)}
    {\(-\frac{4}{5}\)}
\end{problem}

\begin{answer}
D
\end{answer}


\begin{problem}
已知平面 \(\alpha\) 截一球面得圆 \(M\)，过圆心 \(M\) 且与 \(\alpha\) 成 \(60^{\circ}\) 二面角的平面 \(\beta\) 截该球面得圆 \(N\). 若该球面的半径为 \(4\)，圆 \(M\) 的面积为 \(4\pi\)，则圆 \(N\) 的面积为（\quad）

\choices
    {\(7\pi\)}
    {\(9\pi\)}
    {\(11\pi\)}
    {\(13\pi\)}
\end{problem}

\begin{answer}
D
\end{answer}


\begin{problem}
设向量 \(\bs{a}\)，\(\bs{b}\)，\(\bs{c}\) 满足 \(|\bs{a}| = |\bs{b}| = 1\)，\(\bs{a} \cdot \bs{b} = -\frac{1}{2}\)，\(\langle \bs{a} - \bs{c}, \bs{b} - \bs{c} \rangle = 60^\circ\)，则 \(|\bs{c}|\) 的最大值等于（\quad）

\choices
    {\(2\)}
    {\(\sqrt{3}\)}
    {\(\sqrt{2}\)}
    {\(1\)}
\end{problem}

\begin{answer}
A
\end{answer}


\section{填空题}

\begin{problem}
\((1 - \sqrt{x})^{20}\) 的二项展开式中，\(x\) 的系数与 \(x^9\) 的系数之差为\fillinblank{}.
\end{problem}

\begin{answer}
0
\end{answer}


\begin{problem}
已知 \(\alpha \in \left(\frac{\pi}{2}, \pi\right)\)，\(\sin \alpha = \frac{\sqrt{5}}{5}\)，则 \(\tan 2\alpha = \)\fillinblank{}.
\end{problem}

\begin{answer}
\(-\frac{4}{3}\)
\end{answer}


\begin{problem}
已知 \(F_{1}\)，\(F_{2}\) 分别为双曲线 \(C: \frac{x^{2}}{9} - \frac{y^{2}}{27} = 1\) 的左，右焦点，点 \(A \in C\)，点 \(M\) 的坐标为 \((2,0)\)，\(AM\) 为 \(\angle F_{1}AF_{2}\) 的平分线. 则 \(|AF_{2}| = \)\fillinblank{}.
\end{problem}

\begin{answer}
6
\end{answer}


\begin{problem}
已知点 \(E\)，\(F\) 分别在正方体 \(ABCD - A_{1}B_{1}C_{1}D_{1}\) 的棱 \(BB_{1}\)，\(CC_{1}\) 上，且 \(B_{1}E = 2EB\)，\(CF = 2FC_{1}\)，则面 \(AEF\) 与面 \(ABC\) 所成的二面角的正切值等于\fillinblank{}.
\end{problem}

\begin{answer}
\(\frac{\sqrt{2}}{3}\)
\end{answer}


\section{解答题}

\begin{problem}
\(\triangle ABC\) 的内角 \(A\)，\(B\)，\(C\) 的对边分别为 \(a\)，\(b\)，\(c\). 已知 \(A - C = 90^{\circ}\)，\(a + c = \sqrt{2} b\)，求 \(C\).
\end{problem}

\begin{solution}
由三角形内角和及 \(A-C=90^{\circ}\)，得 \(A=C+90^{\circ}\)，\(B=90^{\circ}-2C\)，所以 \(0<C<45^{\circ}\). 根据正弦定理，条件 \(a+c=\sqrt{2}b\) 等价于
\[
\sin A+\sin C=\sqrt{2}\sin B.
\]
代入 \(A=C+90^{\circ}\) 和 \(B=90^{\circ}-2C\)，得 \(\cos C+\sin C=\sqrt{2}\cos 2C\)，即
\[
\cos\left(C-45^{\circ}\right)=\cos 2C.
\]
因为 \(C-45^{\circ}\in(-45^{\circ},0)\)，而 \(2C\in(0,90^{\circ})\)，两角的余弦相等时只能有 \(C-45^{\circ}=-2C\). 因此 \(3C=45^{\circ}\)，故 \(C=15^{\circ}\).
\end{solution}


\begin{problem}
根据以往统计资料，某地车主购买甲种保险的概率为 \(0.5\)，购买乙种保险但不购买甲种保险的概率为 \(0.3\). 设各车主购买保险相互独立.
\begin{enumerate}
\item 求该地 \(1\) 位车主至少购买甲，乙两种保险中的 \(1\) 种的概率；
\item \(X\) 表示该地的 \(100\) 位车主中，甲，乙两种保险都不购买的车主数. 求 \(X\) 的期望.
\end{enumerate}
\end{problem}

\begin{solution}
设事件 \(A\) 为车主购买甲种保险，事件 \(B\) 为车主购买乙种保险. 题设给出 \(P(A)=0.5\)，\(P(B\cap A^{\mathrm c})=0.3\).

\begin{enumerate}
\item 事件“至少购买一种保险”可以分解为互不相容的事件 \(A\) 和 \(B\cap A^{\mathrm c}\)，所以所求概率为 \(P(A)+P(B\cap A^{\mathrm c})=0.5+0.3=0.8\).

\item 一位车主两种保险都不购买的概率为
\[
P(A^{\mathrm c}\cap B^{\mathrm c})=1-P(A)-P(B\cap A^{\mathrm c})=1-0.5-0.3=0.2.
\]
由于各车主购买保险相互独立，\(X\) 服从参数为 \(100\)，\(0.2\) 的二项分布，因此
\[
E(X)=100\times0.2=20.
\]
\end{enumerate}
\end{solution}


\begin{problem}
如图，四棱锥 \(S - ABCD\) 中，\(AB \parallel CD\)，\(BC \perp CD\)，侧面 \(SAB\) 为等边三角形. \(AB = BC = 2\)，\(CD = SD = 1\).
\begin{enumerate}
\item 证明： \(SD \perp\) 平面 \(SAB\)；
\item 求 \(AB\) 与平面 \(SBC\) 所成角的正弦值.
\end{enumerate}
\bitmapfigure[max width=0.55\linewidth]{img/2011/syllabus_science/q19_fig1.png}
\end{problem}

\begin{solution}
取底面 \(ABCD\) 所在平面为 \(z=0\)，建立直角坐标系，使得
\(D=(0,0,0)\)，\(C=(1,0,0)\)，\(B=(1,2,0)\)，\(A=(-1,2,0)\).
设 \(S=(u,v,w)\)，由侧面 \(SAB\) 为边长为 \(2\) 的等边三角形，得 \(SA=SB=2\). 于是
\[
(u+1)^2+(v-2)^2+w^2=(u-1)^2+(v-2)^2+w^2,
\]
从而 \(u=0\). 又因为 \(SD=1\)，\(SA=2\)，所以
\[
\begin{cases}
v^2+w^2=1,\\
1+(v-2)^2+w^2=4.
\end{cases}
\]
两式相减得 \(v=\frac12\)，再由 \(v^2+w^2=1\) 且取 \(w>0\)，得 \(w=\frac{\sqrt{3}}2\). 因而 \(S=\left(0,\frac12,\frac{\sqrt{3}}2\right)\).

\begin{enumerate}
\item 有 \(\overrightarrow{DS}=\left(0,\frac12,\frac{\sqrt{3}}2\right)\)，\(\overrightarrow{AB}=(2,0,0)\)，\(\overrightarrow{AS}=\left(1,-\frac32,\frac{\sqrt{3}}2\right)\). 直接计算得 \(\overrightarrow{DS}\cdot\overrightarrow{AB}=0\)，且 \(\overrightarrow{DS}\cdot\overrightarrow{AS}=-\frac34+\frac34=0\). 由于 \(AB\) 与 \(AS\) 是平面 \(SAB\) 内相交的两条直线，所以 \(DS\perp\) 平面 \(SAB\)，即 \(SD\perp\) 平面 \(SAB\).

\item 取 \(\overrightarrow{SB}=\left(1,\frac32,-\frac{\sqrt{3}}2\right)\)，\(\overrightarrow{SC}=\left(1,-\frac12,-\frac{\sqrt{3}}2\right)\)，则平面 \(SBC\) 的一个法向量为
\[
\bs{n}=\overrightarrow{SB}\times\overrightarrow{SC}=(-\sqrt{3},0,-2).
\]
设 \(AB\) 与平面 \(SBC\) 所成的角为 \(\theta\)，则
\[
\sin\theta=\frac{|\overrightarrow{AB}\cdot\bs{n}|}{|\overrightarrow{AB}|\,|\bs{n}|}=\frac{2\sqrt{3}}{2\sqrt{7}}=\frac{\sqrt{21}}7.
\]
\end{enumerate}
\end{solution}


\begin{problem}
设数列 \(\{a_{n}\}\) 满足 \(a_1 = 0\) 且 \(\frac{1}{1 - a_{n + 1}} -\frac{1}{1 - a_n} = 1\).
\begin{enumerate}
\item 求 \( \{a_{n}\} \) 的通项公式；
\item 设 \( b_n = \frac{1 - \sqrt{a_{n+1}}}{\sqrt{n}} \)，记 \( S_n = \sum_{k=1}^{n} b_k \)，证明：\( S_n < 1 \).
\end{enumerate}
\end{problem}

\begin{solution}
令 \(c_n=\frac{1}{1-a_n}\)，则 \(c_1=1\)，且题设递推式给出 \(c_{n+1}-c_n=1\). 因此 \(\{c_n\}\) 是首项为 \(1\)、公差为 \(1\) 的等差数列，故 \(c_n=n\)，从而
\[
a_n=1-\frac1n=\frac{n-1}{n}.
\]

由此 \(a_{n+1}=\frac{n}{n+1}\)，所以
\[
b_n=\frac{1-\sqrt{\frac{n}{n+1}}}{\sqrt n}=\frac{\sqrt{n+1}-\sqrt n}{\sqrt n\sqrt{n+1}}=\frac1{\sqrt n}-\frac1{\sqrt{n+1}}.
\]
于是
\[
S_n=\sum_{k=1}^n\left(\frac1{\sqrt k}-\frac1{\sqrt{k+1}}\right)=1-\frac1{\sqrt{n+1}}<1.
\]
\end{solution}


\begin{problem}
已知 \(O\) 为坐标原点，\(F\) 为椭圆 \(C: x^2 + \frac{y^2}{2} = 1\) 在 \(y\) 轴正半轴上的焦点，过 \(F\) 且斜率为 \(-\sqrt{2}\) 的直线 \(l\) 与 \(C\) 交于 \(A\)，\(B\) 两点，点 \(P\) 满足 \(\overrightarrow{OA} + \overrightarrow{OB} + \overrightarrow{OP} = 0\).
\begin{enumerate}
\item 证明：点 \(P\) 在 \(C\) 上；
\item 设点 \(P\) 关于点 \(O\) 的对称点为 \(Q\)，证明： \(A\)，\(P\)，\(B\)，\(Q\) 四点在同一圆上.
\end{enumerate}
\FigureLayoutDeclare{img/2011/syllabus_science/q21_fig1.png}{5.5cm}{4.4bp 4.4bp 4bp 4bp}
\bitmapfigure[max width=0.34\linewidth]{img/2011/syllabus_science/q21_fig1.png}
\end{problem}

\begin{solution}
椭圆 \(C\) 的长半轴平方为 \(2\)，短半轴平方为 \(1\)，所以焦距的一半为 \(\sqrt{2-1}=1\)，且 \(F=(0,1)\). 直线 \(l\) 的方程为 \(y=1-\sqrt{2}x\).

\begin{enumerate}
\item 设 \(A=(x_1,y_1)\)，\(B=(x_2,y_2)\). 将直线方程代入椭圆方程，得
\[
x^2+\frac{(1-\sqrt{2}x)^2}{2}=1\quad\Longleftrightarrow\quad 4x^2-2\sqrt{2}x-1=0.
\]
由韦达定理，\(x_1+x_2=\frac{\sqrt{2}}2\). 又 \(y_i=1-\sqrt{2}x_i\)，所以 \(y_1+y_2=1\). 根据向量关系，点 \(P\) 的坐标为
\[
P=\left(-(x_1+x_2),-(y_1+y_2)\right)=\left(-\frac{\sqrt{2}}2,-1\right).
\]
代入椭圆方程，得 \(\left(-\frac{\sqrt{2}}2\right)^2+\frac{(-1)^2}{2}=\frac12+\frac12=1\)，故点 \(P\) 在 \(C\) 上.

\item 由前一项知点 \(Q=\left(\frac{\sqrt{2}}2,1\right)\). 考虑圆
\[
\Gamma:\quad x^2+y^2+\frac{\sqrt{2}}4x-\frac14y-\frac32=0.
\]
其圆心为 \(\left(-\frac{\sqrt{2}}8,\frac18\right)\)，半径平方为 \(\frac{99}{64}>0\)，所以确为一个圆. 对于直线 \(l\) 上的点 \((x,1-\sqrt{2}x)\)，圆 \(\Gamma\) 的左端化为
\[
3x^2-\frac{3\sqrt{2}}2x-\frac34=\frac34\left(4x^2-2\sqrt{2}x-1\right).
\]
因此交点 \(A\)，\(B\) 均在圆 \(\Gamma\) 上. 对点 \(P\) 和 \(Q\) 分别代入圆方程，线性项分别为 \(-\frac14+\frac14=0\) 和 \(\frac14-\frac14=0\)，且 \(x^2+y^2=\frac32\)，故 \(P\)，\(Q\) 也在圆 \(\Gamma\) 上. 所以 \(A\)，\(P\)，\(B\)，\(Q\) 四点在同一圆上.
\end{enumerate}
\end{solution}


\section{选考题：解答题}

\begin{problem}[A]
设函数 \( f(x) = \ln (1 + x) - \frac{2x}{x + 2} \). 证明：当 \( x > 0 \) 时，\( f(x) > 0 \).
\end{problem}

\begin{solution}
函数的定义域为 \(x>-1\)，且 \(f(0)=0\). 对 \(x>-1\) 求导，得
\[
f'(x)=\frac1{1+x}-\frac4{(x+2)^2}=\frac{x^2}{(x+1)(x+2)^2}.
\]
当 \(x>0\) 时，分子 \(x^2>0\)，分母 \((x+1)(x+2)^2>0\)，所以 \(f'(x)>0\). 因而 \(f(x)\) 在 \((0,+\infty)\) 上单调递增，故对任意 \(x>0\)，有 \(f(x)>f(0)=0\).
\end{solution}

\reuseproblemnumber
\begin{problem}[B]
从编号 \(1\) 到 \(100\) 的 \(100\) 张卡片中每次随机抽取一张，然后放回，用这种方式连续抽取 \(20\) 次. 设抽得的 \(20\) 个号码互不相同的概率为 \(p\). 证明：\(p < \left(\frac{9}{10}\right)^{19} < \frac{1}{\e^2}\).
\end{problem}

\begin{solution}
连续抽取且每次放回时，第一次抽取没有限制，后面第 \(k+1\) 次抽取要避开前面已经抽到的 \(k\) 个号码，因此
\[
p=\frac{100}{100}\cdot\frac{99}{100}\cdot\frac{98}{100}\cdots\frac{81}{100}=\prod_{k=1}^{19}\left(1-\frac{k}{100}\right).
\]
对 \(k=1,2,\ldots,9\)，将第 \(k\) 个因子和第 \(20-k\) 个因子配对，有
\[
\left(1-\frac{k}{100}\right)\left(1-\frac{20-k}{100}\right)
=\left(\frac9{10}+\frac{10-k}{100}\right)\left(\frac9{10}-\frac{10-k}{100}\right)<\left(\frac9{10}\right)^2.
\]
中间因子为 \(1-\frac{10}{100}=\frac9{10}\)，所以 \(p<\left(\frac9{10}\right)^{19}\).

由（A）已证结论，取 \(x=\frac19\)，得
\[
\ln\frac{10}{9}>\frac{2\cdot\frac19}{2+\frac19}=\frac2{19}.
\]
于是 \(19\ln\frac9{10}=-19\ln\frac{10}{9}<-2\). 由于指数函数单调递增，得到 \(\left(\frac9{10}\right)^{19}<\e^{-2}=\frac1{\e^2}\). 综上，\(p < \left(\frac9{10}\right)^{19} < \frac1{\e^2}\).
\end{solution}
